\documentclass[journal]{IEEEtran}

\usepackage{booktabs}
\usepackage{array}
\usepackage{tabularx}
\usepackage{ragged2e}
\usepackage{hyperref}
\usepackage{url}
\usepackage{microtype}

\hypersetup{
  colorlinks=true,
  linkcolor=black,
  citecolor=black,
  urlcolor=blue
}

\title{A Comparison Study of Infrastructure-Independent Communication
Approaches for Rural Emergency Messaging}

\author{
Tanishq Mudaliar, Hrshita Balakrishnan, Saivel Konar, Selvanavya Thevar\\[0.5ex]
\textit{SIES Graduate School of Technology, Navi Mumbai, India}
}

\begin{document}
\maketitle

\begin{abstract}
Communication during a flood, fire, medical incident, power outage, or
cellular-network failure is difficult when users are distributed across rural
or peri-rural areas. This paper reviews and compares communication approaches
that may support short emergency messages without assuming continuous Internet
access. The study contrasts conventional cellular and Internet messaging,
Wi-Fi/Bluetooth ad-hoc links, LoRaWAN, general-purpose LoRa mesh systems, and
the proposed LoRaResQ workflow. The comparison uses infrastructure dependence,
network topology, payload suitability, range and energy characteristics,
deployment complexity, and operational and security support as evaluation
criteria.
The literature indicates that LoRa is attractive for small, delay-tolerant
messages because of its long range and low power requirements, but also has
important limitations: low throughput, shared-channel contention, duty-cycle
constraints, and variable propagation. Existing projects demonstrate the
feasibility of low-power LoRa communication, but leave a practical gap
between a flexible technical mesh and a simple, alert-first rural operating
workflow. LoRaResQ addresses that gap as a proposed, low-cost system that
connects a phone to a personal LoRa node and allows eligible packets to be
relayed by nearby participant nodes. The paper does not claim hardware
validation; the proposed evaluation plan identifies the evidence required
before deployment or safety-critical use.
\end{abstract}

\begin{IEEEkeywords}
LoRa, LoRaWAN, mesh networking, emergency communication, disaster resilience,
rural connectivity, off-grid communication.
\end{IEEEkeywords}

\section{Introduction}
Mobile phones are powerful emergency tools only when a working device,
electricity supply, radio link, cellular tower, backhaul, and subscribed
service are available. Rural roads, farms, forest edges, shelters, and
disaster-affected localities may not satisfy all of these conditions. A
communication fallback does not need to carry video or general Internet
traffic to be useful. It may only need to communicate a bounded message such
as ``medical help needed'', ``road blocked'', ``I am safe'', or ``assemble at
the school''.

LoRa is a long-range, low-power physical-layer technology designed for small
packets. LoRaWAN is a network protocol and system architecture built around
star-of-stars gateways, while a LoRa mesh uses peer nodes that can forward
packets. These approaches should not be treated as interchangeable. Adelantado
et al.\ show that LoRaWAN performance is constrained by capacity, collisions,
and scalability assumptions \cite{adelantado2017}. Augustin et al.\ similarly
describe the range, data-rate, and energy trade-offs of LoRa-based networks
\cite{augustin2016}. These limitations matter more in an emergency because an
alert should be useful without creating the false impression of guaranteed
delivery.

This paper presents a focused comparison for the design of LoRaResQ, a
proposed community-owned emergency messaging mesh. The contribution is not a
new radio modulation or a claim of superiority over existing platforms.
Instead, it is a requirements-led comparison and a clearly bounded system
proposal that can be implemented and tested with a small number of nodes.

\section{Review Method}
The study searched Google Scholar, publisher pages, arXiv, and official
standards pages in October 2026 using combinations of ``LoRa'', ``LoRaWAN'',
``mesh'', ``off-grid'', ``emergency communication'', ``disaster'', ``rural'',
``scalability'', and ``Meshtastic''. Sources were included when they were
journal articles, scholarly preprints, or primary technical specifications
with a clear description of architecture, constraints, or measured results.
The six scholarly papers selected for this focused comparison are mostly
adjacent LoRa, LoRaWAN, mesh, and Meshtastic studies; LoRaMoto is an
emergency-specific simulation study rather than a field trial. This limitation
is recorded explicitly below. Sources were also
screened for direct relevance to short-message communication rather than
general high-bandwidth networking. Official technical documentation is used
only to clarify terminology, regulation, and protocol behaviour; it is not
treated as independent experimental evidence. Conference proceedings were not
included in this compact review corpus so that the comparison would use
sources with complete, inspectable methods and conclusions; this choice may
exclude relevant emergency-network results.

The comparison criteria are:
\begin{enumerate}
  \item \textbf{Infrastructure dependence:} whether the system requires a
  cellular network, Internet, gateway, or cloud service.
  \item \textbf{Network topology:} direct, star, gateway-based, or multi-hop
  peer communication.
  \item \textbf{Payload suitability:} whether the system is appropriate for
  short text and structured alerts rather than high-bandwidth content.
  \item \textbf{Range and energy:} reported or expected suitability for
  battery-powered portable nodes, without assuming a universal range.
  \item \textbf{Deployment complexity:} hardware, gateway, configuration, and
  operational requirements.
  \item \textbf{Operational and security support:} delivery states,
  prioritisation, expiry, localisation, authentication, and privacy.
\end{enumerate}

\section{Technologies and Related Systems}
\subsection{Conventional connected communication}
Cellular data, SMS, and Internet messaging provide familiar user experiences
and can support large payloads. Their weakness in this context is dependence on
external infrastructure. A damaged tower, congested cell, unavailable
backhaul, depleted battery, or missing coverage can prevent a message from
being sent. These services should therefore remain the preferred channel when
available, while an off-grid system should be treated as a fallback rather
than a replacement for official emergency services. This paper treats
cellular and Internet messaging as the baseline against which the
infrastructure dependence of the proposed system is compared; it does not
claim that the baseline is technically inferior.

\subsection{Wi-Fi and Bluetooth ad-hoc links}
Wi-Fi and Bluetooth Low Energy are useful for connecting a phone to a nearby
device or for local peer-to-peer communication. Bluetooth Low Energy is
specified for low-power short-range links \cite{bluetoothcore}; Wi-Fi provides
higher throughput than LoRa under the IEEE 802.11 family of local-area
network standards \cite{ieee802}, but Wi-Fi and Bluetooth links have shorter
practical range and greater sensitivity to obstructions and user density. A
hybrid architecture can use Bluetooth for
the phone-to-node link and LoRa for the outdoor node-to-node link. This
separates the user interface from the low-bandwidth radio network.

\subsection{LoRaWAN}
LoRaWAN provides a standardised low-power wide-area architecture, normally
using end devices, gateways, and network/application servers. Its strengths
include mature device classes, long-range links, and low energy consumption.
Its architecture is not inherently an infrastructure-free peer mesh: a
gateway and network service are normally required to carry data to an
application. Recent survey work on scalable LoRaWAN similarly identifies traffic,
interference, resource allocation, and topology as important constraints
\cite{jouhari2023}; the limits identified by Adelantado et al.\ include
collisions, data-rate trade-offs, and duty-cycle-related constraints
\cite{adelantado2017}.
The LoRaWAN specification defines the end-device, gateway, and network-server
architecture used by this comparison \cite{lorawan2020}.

\subsection{LoRa mesh systems}
Peer-to-peer LoRa systems remove the requirement for a permanent gateway by
allowing nodes to relay selected packets. A recent journal review comparing
open-source and proprietary LoRa mesh technologies identifies cost, energy
efficiency, latency, range, security, and routing as the main design
trade-offs \cite{hammouti2024}. A mesh can extend reach when a direct path is
blocked, but every relay consumes airtime and energy. Duplicate suppression,
hop limits, queue policy, packet expiry, and acknowledgement semantics are
therefore essential.
LoRaMoto is a closer emergency-specific example: it extends LoRaWAN with
end-node forwarding for civilian safety-status messages after an earthquake,
but its evaluation is simulation-based rather than a field validation
\cite{loramoto2021}.

\subsection{Meshtastic as a practical reference}
Meshtastic is an important open-source reference for practical off-grid LoRa
mesh communication. Its official documentation describes managed flooding with
hop-limited, SNR-aware rebroadcast, next-hop routing for direct messages since
firmware 2.6, and CSMA/CA channel access \cite{meshtasticdocs}. It
demonstrates that inexpensive nodes, user devices, and mesh forwarding can be
combined in a community project. Its acknowledgement semantics are also
instructive: a broadcast is treated as implicitly acknowledged once the sender
hears any node rebroadcast it, whether or not that node holds the channel key
\cite{meshtasticdocs}. This is the local-acceptance versus
recipient-acknowledgement distinction that LoRaResQ's status model is meant to
expose. A recent Meshtastic-based campus system reports a maximum field link of
approximately 2.47 km under its stated campus conditions \cite{meshtastic2026}.
The paper's body does not establish that distance as a statistically
validated packet-delivery result: its reported tracker frames, duty-cycle
setting, and 22 dBm test configuration require careful interpretation. The
authors identify packet-delivery-rate and latency campaigns as future work.
This arXiv preprint concerns campus telemetry, not emergency response, and
must not be presented as a LoRaResQ range guarantee.
LoRaResQ does not claim to replace or outperform Meshtastic. It narrows the
user experience to a rural emergency workflow: a small set of preset alerts,
bounded text, explicit status, local-language labels, and a deployment plan
that can be demonstrated with three interoperable personal nodes.

\section{Comparison}
Table~\ref{tab:comparison} compares the alternatives against the requirements
of a local emergency-message fallback. The qualitative entries describe
architectural fit, not measured performance.

\begin{table*}[t]
\caption{Qualitative comparison for rural emergency messaging}
\label{tab:comparison}
\centering
\scriptsize
\begin{tabularx}{\textwidth}{@{}>{\raggedright\arraybackslash}X
>{\raggedright\arraybackslash}X
>{\raggedright\arraybackslash}X
>{\raggedright\arraybackslash}X
>{\raggedright\arraybackslash}X
>{\raggedright\arraybackslash}X@{}}
\toprule
\textbf{Approach} & \textbf{Infrastructure} & \textbf{Topology} &
\textbf{Payload} & \textbf{Range / energy} & \textbf{Deployment, operations,
and security} \\
\midrule
Cellular / Internet & Required & Provider infrastructure &
High bandwidth & Phone-dependent & Mature; unavailable during outage \\
Wi-Fi ad hoc & Not necessarily & Direct / configurable mesh &
High bandwidth, short range & Short range; higher power draw & Configuration-heavy; security depends on setup \\
Bluetooth LE & Not required & Direct / local relay &
Small payloads & Good for short links & Good phone link; not wide-area \\
LoRaWAN & Gateway and server normally required & Star-of-stars &
Small, delay-tolerant & Good low-power range & Mature, but gateway and server required \\
Generic LoRa mesh & No permanent gateway required & Multi-hop peer &
Small, delay-tolerant & Good, relay-dependent & Routing and security policy required \\
Meshtastic reference & Configurable; can be off-grid & Multi-hop peer &
Small text and data & Good, topology-dependent & Mature reference; broad configuration surface; encrypted channels available \\
\textbf{LoRaResQ (proposed)} & \textbf{No permanent gateway for prototype} &
\textbf{Bounded multi-hop} & \textbf{Alerts and short text} &
\textbf{Battery-oriented (target)} & \textbf{Alert-first; status model and security planned; unvalidated, no field data} \\
\bottomrule
\end{tabularx}
\end{table*}

The comparison produces three design conclusions. First, no approach removes
the need to state operating conditions: ``sent to the local node'' is not the
same as ``received by the intended person''. Second, LoRa's advantage is
appropriate only when the payload is small and delay is acceptable. Third, the
user workflow is a major part of emergency-system reliability. A system that
has a radio link but gives ambiguous status, stale alerts, or difficult
configuration may fail operationally even if its physical layer works.

\section{Proposed LoRaResQ System}
LoRaResQ is designed as a small, community-owned mesh for hamlets, farms,
forest-edge routes, shelters, or a contained event. Each participant carries
a personal node. The phone communicates with its node over Bluetooth LE; the
node sends a bounded LoRa packet. An eligible new packet may be forwarded by
another participant's node, subject to a hop limit, duplicate cache, queue
policy, and packet expiry. The prototype target is three nodes in the path
\mbox{A $\rightarrow$ B $\rightarrow$ C}, where A and C cannot communicate
directly.

The application is planned to support:
\begin{itemize}
  \item structured alerts such as medical help, fire, flood, road blocked,
  assemble, and all safe;
  \item bounded free-text messages;
  \item states that distinguish queued, accepted by the node, transmitted,
  relayed, and acknowledged;
  \item local-language labels and icon-based actions;
  \item priority and expiry rules to prevent stale or excessive traffic; and
  \item a deployment record containing region, radio configuration, node
  identity, and test observations.
\end{itemize}

Security is part of the proposed design rather than an achieved result. For
one-to-one messages, the nodes should use authenticated encryption with
per-device or per-recipient keys; message identifiers, nonces, hop limits,
and replay protection should prevent duplicated or forged packets. Key
enrolment and revocation must be defined before field use, and emergency
alerts should not contain unnecessary personal or medical data. The prototype
will demonstrate the interface and packet-state model first; cryptographic
interoperability and resistance to impersonation remain validation tasks.
Broadcast alerts need separate protection: because any node can originate a
broadcast, alerts should carry sender authentication, for example a
per-community key or per-node signature, and relaying nodes should forward
one-to-one messages as ciphertext they cannot read.

The system is deliberately not specified as a guaranteed life-safety service.
It cannot dispatch an ambulance automatically, guarantee delivery, replace
112, provide Internet access, or substitute for trained responders. When
cellular or official emergency communication works, users should use it.

\section{India-Specific Regulatory Context}
For a rural Indian prototype, the radio configuration must be selected only
after checking the current Wireless Planning and Coordination (WPC) rules.
The 2021 exemption rules (G.S.R.~853(E), 10 December 2021) allow low-power
short-range devices in the 865--868 MHz band on a non-interference,
non-protection, shared and non-exclusive basis, subject to per-category limits
\cite{india2021}. Two categories are candidates for an alert-messaging node:
non-specific short-range devices (25 mW e.r.p., 1\% duty cycle, with an FHSS
condition on occupied bandwidth) and tracking, tracing and data-acquisition
devices (500 mW e.r.p., adaptive power control, duty cycle $\leq$10\% for
network access points and $\leq$2.5\% otherwise, bandwidth $\leq$200 kHz).
Which category covers a 125 kHz LoRa emergency-messaging node has not been
confirmed with WPC, so the prototype is designed against the stricter limits.
Rule 5 requires covered equipment to conform to these limits, to be type
approved, and to comply with the corresponding EN standard \cite{india2021}.
An EU868 preset must not be assumed legal in India. The transmit power,
bandwidth, and duty-cycle settings used in any trial belong in the deployment
record; they are not results claimed by this literature review.
The 22 dBm campus configuration reported in the Meshtastic study must not be
copied into an India trial: it exceeds the nominal 25 mW limit before
antenna-gain and e.r.p. calculations. Firmware should enforce the selected
power cap rather than relying on a regional preset.

\section{Proposed Evaluation Plan}
Because the project currently has no hardware, all hardware-related values in
this paper are requirements or targets rather than results. The planned
validation should include:
\begin{enumerate}
  \item a controlled two-node test with one hundred short messages, selected
  to provide an initial estimate of delivery proportion rather than a
  statistically generalisable reliability claim;
  \item twenty three-node forwarding attempts as a feasibility smoke test,
  with the direct A-C link unavailable or intentionally attenuated;
  \item outdoor trials in open and obstructed terrain, recording distance,
  antenna, spreading-factor settings, transmit power, RSSI/SNR, airtime and
  duty cycle, packet loss, and acknowledgements;
  \item battery measurements under a declared workload;
  \item duplicate and stale-packet tests; and
  \item phone-to-node Bluetooth Low Energy range, reconnection, and message
  transfer tests for the hybrid interface;
  \item a short usability test measuring whether a first-time user can send
  an alert within thirty seconds without confusing local acceptance with
  recipient acknowledgement. The thirty-second target is an engineering
  usability target for the prototype, not a safety standard.
\end{enumerate}

The results should be reported with the test environment and limitations,
including delivery proportions with a 95\% confidence interval. The two-node
and three-node tests are prototype targets; if three radio modules are not
available, the forwarding test must be reported as emulated rather than as a
three-node hardware result. Vendor range claims should not be presented as
project results. No deployment
should be described as public-safety-ready until radio compliance, electrical
safety, security, and field reliability have been reviewed.

\section{Limitations of the Study}
This is a focused design-oriented literature review, not a systematic review
with a registered search protocol or meta-analysis. Published systems use
different hardware, frequencies, antennas, firmware, terrain, payload sizes,
and traffic loads, so their numerical results cannot be transferred directly
to LoRaResQ. The proposed system has not yet been tested because the team does
not currently possess the radio hardware. The comparison therefore supports
the design rationale and test plan; it does not prove range, battery life,
reliability, or emergency effectiveness. In addition, the scholarly studies
are mostly adjacent technology sources rather than emergency or rural field
trials; LoRaMoto is emergency-specific but simulation-based. Excluding
conference proceedings may omit relevant disaster communication evidence. The
proposed security controls are design
requirements, not a completed security evaluation.

\section{Conclusion}
Cellular and Internet services remain the most capable communication options
when available, but they are not dependable as the only channel in every
rural or disaster scenario. Wi-Fi and Bluetooth are useful for local links,
while LoRa provides a promising low-power long-range link for small messages.
LoRaWAN offers standardisation and mature ecosystem support but normally
depends on gateway infrastructure. Peer LoRa mesh systems reduce that
dependency at the cost of airtime, queueing, forwarding, and operational
complexity.

The literature supports the feasibility of investigating LoRa-based
off-grid messaging, but it also warns against overclaiming. LoRaResQ's
distinct contribution is a bounded, alert-first workflow for a small rural
community, not a new radio technology. Its next credible milestone is a
measured three-node prototype whose interface, delivery states, relay policy,
and limitations can be demonstrated transparently.

\begin{thebibliography}{00}

\bibitem{adelantado2017}
F. Adelantado, X. Vilajosana, P. Tuset-Peir\'o, B. Mart\'inez,
  J. Meli\`a-Segu\'i,
  and T. Watteyne, ``Understanding the limits of LoRaWAN,''
  \emph{IEEE Communications Magazine}, vol. 55, no. 9, pp. 34--40, 2017.
  doi: 10.1109/MCOM.2017.1600613.

\bibitem{augustin2016}
A. Augustin, J. Yi, T. Clausen, and W. M. Townsley, ``A study of LoRa:
  Long range \& low power networks for the Internet of Things,''
  \emph{Sensors}, vol. 16, no. 9, Art. no. 1466, 2016.
  doi: 10.3390/s16091466.

\bibitem{bluetoothcore}
Bluetooth SIG, ``Bluetooth Core Specification, Version 5.4,'' technical
  specification, 2023. [Online]. Available:
  \url{https://www.bluetooth.com/specifications/specs/core-specification-5-4/}.

\bibitem{ieee802}
IEEE, ``IEEE Standard for Information Technology--Telecommunications and
  Information Exchange Between Systems Local and Metropolitan Area Networks--
  Specific Requirements--Part 11: Wireless LAN Medium Access Control (MAC)
  and Physical Layer (PHY) Specifications,'' IEEE Std. 802.11-2020, 2020.
  [Online]. Available: \url{https://standards.ieee.org/standard/802_11-2020.html}.

\bibitem{jouhari2023}
M. Jouhari, N. Saeed, M.-S. Alouini, and E. M. Amhoud, ``A survey on
  scalable LoRaWAN for massive IoT: Recent advances, potentials, and
  challenges,'' \emph{IEEE Communications Surveys \& Tutorials}, vol. 25,
  no. 3, pp. 1841--1876, 2023. doi: 10.1109/COMST.2023.3274934.

\bibitem{lorawan2020}
LoRa Alliance, ``LoRaWAN specification,'' technical specification,
2020. [Online]. Available:
 \url{https://lora-alliance.org/wp-content/uploads/2020/11/lorawan1.0.3.pdf}.

\bibitem{hammouti2024}
M. Hammouti, O. Moussaoui, M. Hassine, and A. Betari, ``Exploring open
  source and proprietary LoRa mesh technologies,'' \emph{Indonesian
  Journal of Electrical Engineering and Computer Science}, vol. 34, no. 2,
  pp. 960--969, 2024. doi: 10.11591/ijeecs.v34.i2.pp960-969.

\bibitem{loramoto2021}
R. Pueyo Centelles \emph{et al.}, ``LoRaMoto: A communication system to
  provide safety awareness among civilians after an earthquake,''
  \emph{Future Generation Computer Systems}, 2021.
  doi: 10.1016/j.future.2020.07.040.

\bibitem{meshtasticdocs}
Meshtastic, ``Meshtastic documentation: Mesh broadcast algorithm and
  configuration,'' online documentation, accessed Oct. 2, 2026. [Online].
  Available: \url{https://meshtastic.org/docs/overview/mesh-algo/}.

\bibitem{meshtastic2026}
R. Garzon Andosilla and J. de Jes\'us Rugeles Uribe, ``A Meshtastic-based
  LoRa mesh system for smart campus applications: From solar-powered
  sensing to containerized data management,'' arXiv preprint
  arXiv:2605.20379, 2026. [Online]. Available:
  \url{https://arxiv.org/abs/2605.20379}.

\bibitem{india2021}
Government of India, ``Use of Low Power Equipment in the Frequency Band
  865--868 MHz for Short Range Devices (Exemption from Licence) Rules, 2021,''
  G.S.R. 853(E), 2021. [Online]. Available:
  \url{https://thc.nic.in/Central\%20Governmental\%20Rules/use\%20of\%20low\%20power\%20Equipment\%20in\%20the\%20frequency\%20band\%20865\%20to\%20868\%20MHz\%20for\%20Short\%20Range\%20Devices\%20Exemption\%20from\%20Licence\%20Rules,2021.pdf}.

\end{thebibliography}

\end{document}
